\section{尤度と線形モデル}
\label{b24_linear_model}

\subsection{授業内容}

\subsubsection{概要}


\subsubsection{コマ主題細目}
\begin{description}
	\item[回帰モデルの確率表現] 単回帰モデル$Y_i = \beta_0 + \beta_1 X_i + \varepsilon_i$を確率モデルとして表現し，最尤推定することを考える。正規分布の平均パラメータを線形予測子$\beta_0 + \beta_1 X_i$で表現し，誤差項$\varepsilon_i$が正規分布$N(0, \sigma^2)$に従うと仮定する。このとき，観測データ$Y_i$は平均$\beta_0 + \beta_1 X_i$，分散$\sigma^2$の正規分布に従う確率変数として表現できることを理解する。
    \item[回帰分析の最尤推定] 回帰モデルのパラメータ（切片$\beta_0$，傾き$\beta_1$，誤差分散$\sigma^2$）を最尤推定する方法を学ぶ。尤度関数の構築とその最大化により，最尤推定値を導出する手順を理解する。また，Rによる実践を通じて，最尤推定の結果と最小二乗法による推定結果が一致することを確認する。
    \item[線形モデルの拡張] 線形モデルを一般化し，様々なタイプのデータに対応する方法を学ぶ。線形結合でモデルを表現すること，不確定な誤差の記述として確率分布を用い，確率分布の平均的パラメータに線形予測子を対応させることを理解する。回帰分析はこの枠組みの一つであり，かつ，リンク関数が恒常関数である特別な場合であることを学ぶ。
    \item[一般化線形モデル] 具体的な確率分布を導入し，パラメータの定義域に応じたリンク関数を用いることで，二項分布やポアソン分布などの非正規分布を扱う一般化線形モデル（GLM: Generalized Linear Model）の基礎を理解する。これにより，様々なタイプのデータに対して回帰分析を適用できることを学ぶ。
\end{description}
\subsubsection{キーワード}
\begin{itemize}
	\item 回帰モデルの確率表現
	\item 最尤推定
	\item 線形予測子
	\item リンク関数
	\item 一般化線形モデル
\end{itemize}
\subsection{授業情報}
\paragraph{コマの展開方法}
講義
\paragraph{標準シラバスにおける位置づけ}
科目番号5；心理学統計法；(1)心理学で用いられる統計手法；J より高度な記述や推定を目指して
\subsubsection{予習・復習課題}
\paragraph{予習}
正規分布の確率密度関数を確認し，正規分布を用いた最尤推定の基礎を復習しておくこと。特に尤度と確率の違いについて整理しておくことが肝要である。また回帰分析の基本である，誤差項が正規分布に従うという仮定の根拠，ひいてはデータの背後にある確率モデルの考え方についても再確認しておくとよい。
\paragraph{復習}
  線形予測子とリンク関数の役割を整理し、一般化線形モデルがどのようにして様々なデータ型に対応できるか，自分で説明できるようになることが望ましい。また、最尤推定の手順を再確認し、Rを用いて回帰モデルの最尤推定を実践できるようにしておくこと。具体的には、Rの`glm()`関数を使用して一般化線形モデルを適用し、結果の解釈ができるようになることが目標である。さらに、最小二乗法と最尤法の関係についても理解を深めておくとよい。